%% toc: P15 163 163.10
% Proofs for the comparison of exchangeability with the inductive task.
\section{Exchangeability assumptions and induction: proofs}
\label{sec:annex-p15-w02b-p15-022}

Throughout, \(\Omega=\{0,1\}^{\mathbb N_+}\) or \(\mathbb R^{\mathbb N_+}\)
carries its product sigma-field and coordinates \(X_j\), as stated in each
result, and the notation of \cref{def:p15-w02b-p15-022-task} is used.

\begin{annexitem}{annex:p15-w02b-p15-022-symmetry-consequences}{thm:p15-w02b-p15-022-symmetry-consequences}{What exchangeability commits a law to}
\begin{proof}
Exchangeability means that
\((X_{\pi(1)},\ldots,X_{\pi(n)})\) has the law of \((X_1,\ldots,X_n)\) for
every permutation \(\pi\) of \(\{1,\ldots,n\}\). Equal joint laws have
equal laws of their first coordinate and of their first two coordinates.

(1) Let \(\pi\) exchange \(1\) and \(j\). The first coordinate of the
permuted vector is \(X_j\), so \(X_j\) has the law of \(X_1\).

(2) Since \(i\neq j\) and \(n\geq2\), there is a permutation with
\(\pi(1)=i\) and \(\pi(2)=j\). The first two coordinates of the permuted
vector are \((X_i,X_j)\), which therefore has the law of \((X_1,X_2)\).
The case \(i=2\), \(j=1\) gives the second assertion.

(3) The sequence takes values in \([0,1]\), so
\(E(X_1\mid\mathcal E)\) has a version with values in \([0,1]\); this is
\(\Theta\). By \cref{thm:p15-exchangeable-strong-laws} with
\(f(x)=x\), \(S_n/n\to\Theta\) almost surely. By
\cref{thm:p15-binary-definetti}, the mixing measure \(\lambda\) is the law
of the almost-sure limiting frequency, hence the law of \(\Theta\). The
word of \(k\) ones has count \(k\), so the same theorem gives
\[
 P(X_1=\cdots=X_k=1)=\int_0^1\theta^k\,\lambda(d\theta)=E\Theta^k .
\]
The cases \(k=1,2\) are the displayed identities.
\end{proof}
\end{annexitem}

\begin{annexitem}{annex:p15-w02b-p15-022-predictive}{thm:p15-w02b-p15-022-predictive}{Predictive updating and limiting frequency}
\begin{proof}
As in the preceding proof, \(\Theta\) is a version of
\(E(X_1\mid\mathcal E)\) with values in \([0,1]\), its law is \(\lambda\),
and \(S_n/n\to\Theta\) almost surely. Write
\(\mathcal K=\sigma(\mathcal F_n\cup\mathcal E)\).

(1) Fix \(N\geq n\geq0\). Apply \cref{thm:p15-definetti-factorization}
with \(k=N+1\) and the bounded functions
\(f_j(x)=\mathbf1_{\{x=b_j\}}\) for \(j\leq n\), \(f_j\equiv1\) for
\(n<j\leq N\), and \(f_{N+1}(x)=x\), where \(b\in\{0,1\}^n\) has \(s\)
ones. Since \(E(\mathbf1_{\{X_1=1\}}\mid\mathcal E)=\Theta\),
\(E(\mathbf1_{\{X_1=0\}}\mid\mathcal E)=1-\Theta\) and
\(E(1\mid\mathcal E)=1\) almost surely,
\[
 E\bigl(\mathbf1_{\{X_{1:n}=b\}}X_{N+1}\mid\mathcal E\bigr)
       =\Theta^{s+1}(1-\Theta)^{n-s}.
\]
The same theorem with \(k=n\) gives
\(E(\mathbf1_{\{X_{1:n}=b\}}\mid\mathcal E)=\Theta^s(1-\Theta)^{n-s}\);
for \(n=0\) both identities hold with the empty product equal to one.
For \(C\in\mathcal E\), multiply by \(\mathbf1_C\), which is
\(\mathcal E\)-measurable and bounded, and integrate:
\[
 \begin{aligned}
 E\bigl(X_{N+1};\{X_{1:n}=b\}\cap C\bigr)
   &=E\bigl(\mathbf1_C\Theta^{s+1}(1-\Theta)^{n-s}\bigr)\\
   &=E\bigl(\Theta\mathbf1_C\,E(\mathbf1_{\{X_{1:n}=b\}}\mid\mathcal E)\bigr)
   =E\bigl(\Theta;\{X_{1:n}=b\}\cap C\bigr).
 \end{aligned}
\]
The third equality uses that \(\Theta\mathbf1_C\) is bounded and
\(\mathcal E\)-measurable. The sets \(\{X_{1:n}=b\}\cap C\), together
with \(\varnothing\), form a pi-system: two of them intersect in
\(\{X_{1:n}=b\}\cap(C\cap C')\) when their words agree and in
\(\varnothing\) otherwise. It generates \(\mathcal K\), since
\(\{X_{1:n}=b\}=\{X_{1:n}=b\}\cap\Omega\) and
\(C=\bigcup_b(\{X_{1:n}=b\}\cap C)\). The class of
\(A\in\mathcal K\) with \(E(X_{N+1};A)=E(\Theta;A)\) contains
\(\Omega\), the finite disjoint union of the sets
\(\{X_{1:n}=b\}\cap\Omega\); it is closed under proper differences by
linearity and under increasing unions by dominated convergence, both
integrands lying in \([0,1]\). By the pi--lambda theorem it is all of
\(\mathcal K\). Since \(\Theta\) is \(\mathcal K\)-measurable, it is a
version of \(E(X_{N+1}\mid\mathcal K)\). As
\(\mathcal F_n\subseteq\mathcal K\), the tower property gives
\(P(X_{N+1}=1\mid\mathcal F_n)=E(X_{N+1}\mid\mathcal F_n)
=E(\Theta\mid\mathcal F_n)\), and the right side does not depend on \(N\).

(2) The word \(b1\) of length \(n+1\) has \(s+1\) ones. By
\cref{thm:p15-binary-definetti},
\[
 \begin{aligned}
 P(X_{1:n}=b,\ X_{n+1}=1)&=\int\theta^{s+1}(1-\theta)^{n-s}\,\lambda(d\theta),\\
 P(X_{1:n}=b)&=\int\theta^{s}(1-\theta)^{n-s}\,\lambda(d\theta).
 \end{aligned}
\]
Their ratio is the displayed formula; both integrals depend on \(b\) only
through \(n\) and \(s\).

(3) By (1), \(Q_n=E(\Theta\mid\mathcal F_n)\) almost surely. The fields
\(\mathcal F_n\) increase, and \(\sigma(\bigcup_n\mathcal F_n)=\mathcal F\)
because the coordinates generate the product sigma-field. \(\Theta\) is
bounded and \(\mathcal E\subseteq\mathcal F\), so it is
\(\mathcal F\)-measurable. \Cref{thm:p15-levy-upward} gives
\(Q_n\to E(\Theta\mid\mathcal F)=\Theta\) almost surely and in \(L^1\).
Together with \(S_n/n\to\Theta\) almost surely, the difference tends to
zero almost surely.
\end{proof}
\end{annexitem}

\begin{annexitem}{annex:p15-w02b-p15-022-mixing-law}{thm:p15-w02b-p15-022-mixing-law}{The same symmetry with different mixing laws}
\begin{proof}
In each case \cref{thm:p15-w02b-p15-022-predictive}, item~2, applies with
\(n\) and \(s=s_n\).

(1) For \(\lambda=\delta_p\) the word has probability
\(p^{s}(1-p)^{n-s}\) and the numerator is \(p\) times it, so the ratio is
\(p\) whenever the denominator is positive.

(2) For \(a,b>0\) put \(B(a,b)=\int_0^1\theta^{a-1}(1-\theta)^{b-1}\,d\theta\),
finite and positive. We show \(B(a+1,b)=aB(a,b)/(a+b)\). On
\([\varepsilon,1-\varepsilon]\), integration by parts with
\(u=(1-\theta)^b\) and \(v=\theta^a/a\) gives
\[
 \int_\varepsilon^{1-\varepsilon}\theta^{a-1}(1-\theta)^b\,d\theta
 =\Bigl[\frac{\theta^a}{a}(1-\theta)^b\Bigr]_\varepsilon^{1-\varepsilon}
 +\frac ba\int_\varepsilon^{1-\varepsilon}\theta^a(1-\theta)^{b-1}\,d\theta .
\]
As \(\varepsilon\downarrow0\) the bracket tends to zero because \(a,b>0\),
and both integrals increase to their limits by monotone convergence. Hence
\(B(a,b+1)=(b/a)B(a+1,b)\). Also
\(\theta^{a-1}(1-\theta)^{b-1}=\theta^a(1-\theta)^{b-1}+\theta^{a-1}(1-\theta)^b\)
gives \(B(a,b)=B(a+1,b)+B(a,b+1)=B(a+1,b)(1+b/a)\), as claimed.

For \(\lambda=\operatorname{beta}(\alpha,\beta)\), with density
\(\theta^{\alpha-1}(1-\theta)^{\beta-1}/B(\alpha,\beta)\), put
\(a=\alpha+s\), \(b=\beta+n-s\). The word has probability
\(B(a,b)/B(\alpha,\beta)>0\), and the ratio in item~2 is
\(B(a+1,b)/B(a,b)=a/(a+b)=(\alpha+s)/(\alpha+\beta+n)\).
The uniform law is \(\alpha=\beta=1\), giving
\[
 \frac{s_n+1}{n+2}=\frac{s_n/n+1/n}{1+2/n}\longrightarrow f .
\]

(3) Write \(u=(1/10)^{s}(9/10)^{n-s}\) and \(v=(9/10)^s(1/10)^{n-s}\).
The word has probability \((u+v)/2>0\), and the numerator is
\((u/10+9v/10)/2\). With \(R=v/u=9^{2s-n}\),
\[
 p_n(\lambda)=\frac{1/10+(9/10)R}{1+R}.
\]
Since \(2s_n-n=n(2s_n/n-1)\), \(R\to\infty\) when \(f>1/2\) and
\(R\to0\) when \(f<1/2\). The two limits are \(9/10\) and \(1/10\).

(4) Let \(U\) be the union of all relatively open subsets of \([0,1]\) of
\(\lambda\)-measure zero; the support is \([0,1]\setminus U\). Each point of
\(U\) lies in a relatively open interval with rational endpoints contained in
a \(\lambda\)-null open set, so \(U\) is a countable union of null sets and
\(\lambda(U)=0\). As \(\Theta\) has law \(\lambda\), \(P_\lambda(\Theta\in
U)=0\). The set of sequences whose frequency converges to a point of \(U\)
is contained in \(\{\Theta\in U\}\cup\{S_n/n\not\to\Theta\}\), which is
null by the first part and \cref{thm:p15-w02b-p15-022-predictive}.
\end{proof}
\end{annexitem}

\begin{annexitem}{annex:p15-w02b-p15-022-recoding}{thm:p15-w02b-p15-022-recoding}{Exchangeability depends on the description}
\begin{proof}
(1) For \(j\leq N\), \(Y_j=X_j\) by definition.

(2) If the \(X_j\) are independent fair variables, the \(Y_j\) are
independent, being functions of distinct independent coordinates, and each
is fair. Their finite laws give every word of length \(n\) probability
\(2^{-n}\), which is invariant under rearrangement; so \(Y\) is
exchangeable.

Conversely suppose \(Y\) exchangeable. Put \(m=P(X_1=1)\). By
\cref{thm:p15-w02b-p15-022-symmetry-consequences}, item~1, applied to
\(X\), every \(P(X_j=1)=m\). Since \(N\geq1\), \(P(Y_1=1)=m\) and
\(P(Y_{N+1}=1)=P(X_{N+1}=0)=1-m\); item~1 applied to \(Y\) gives
\(m=1/2\). By item~2 for \(X\), every pair of distinct indices has
\(P(X_i=1,X_j=1)=a\) for one number \(a\). Then
\[
 P(X_i=0,X_j=0)=1-2m+a=a,\qquad P(X_i=1,X_j=0)=m-a=\tfrac12-a .
\]
For \(Y\), the distinct pairs \((1,N+1)\) and \((N+1,N+2)\) give
\[
 \begin{aligned}
 P(Y_1=1,Y_{N+1}=1)&=P(X_1=1,X_{N+1}=0)=\tfrac12-a,\\
 P(Y_{N+1}=1,Y_{N+2}=1)&=P(X_{N+1}=0,X_{N+2}=0)=a .
 \end{aligned}
\]
Item~2 applied to \(Y\) makes these equal, so \(a=1/4\). By item~3,
\(E\Theta=1/2\) and \(E\Theta^2=1/4\), so
\(\operatorname{Var}\Theta=0\) and \(\Theta=1/2\) almost surely. Then
\(\lambda=\delta_{1/2}\), and \cref{thm:p15-binary-definetti} gives every
word of length \(n\) probability \(2^{-n}\). These are the finite laws of
independent fair coordinates; since finite-dimensional laws determine a law
on the product sigma-field, \(X\) is independent and fair.

(3) Let \(\rho(x)=(x_1,\ldots,x_N,1-x_{N+1},1-x_{N+2},\ldots)\). It is a
measurable involution of \(\Omega\) and \(Y=\rho(X)\). Put
\(P'=P_{\mathrm{unif}}\circ\rho^{-1}\). Under \(P'\) the law of
\(Y=\rho(X)\) is \(P_{\mathrm{unif}}\circ\rho^{-1}\circ\rho^{-1}=P_{\mathrm{unif}}\),
so \(Y\) is exchangeable with uniform mixing law and such a \(P'\) exists.
Conversely, for any \(P'\) under which \(Y\) is exchangeable with uniform
mixing law, \cref{thm:p15-binary-definetti} fixes every finite law of \(Y\),
so the law of \(Y\) under \(P'\) is \(P_{\mathrm{unif}}\). For
\(B\subseteq\{0,1\}^N\),
\[
 P'(X_{1:N}\in B)=P'(Y_{1:N}\in B)=P_{\mathrm{unif}}(X_{1:N}\in B),
\]
so the two laws agree on every event determined by the first \(N\)
instances. By \cref{thm:p15-w02b-p15-022-mixing-law}, item~2, with
\(s=n=N\), \(P_{\mathrm{unif}}(X_{N+1}=1\mid X_{1:N}=1\cdots1)=(N+1)/(N+2)\).
The same theorem applied to \(Y\) under \(P'\) gives
\(P'(Y_{N+1}=1\mid Y_{1:N}=1\cdots1)=(N+1)/(N+2)\). Since
\(\{Y_{1:N}=1\cdots1\}=\{X_{1:N}=1\cdots1\}\) and
\(\{X_{N+1}=1\}=\{Y_{N+1}=0\}\),
\(P'(X_{N+1}=1\mid X_{1:N}=1\cdots1)=1-(N+1)/(N+2)=1/(N+2)\).
\end{proof}
\end{annexitem}

\begin{annexitem}{annex:p15-w02b-p15-022-markov-trend}{thm:p15-w02b-p15-022-markov-trend}{A stationary chain and a trend violate the symmetry}
\begin{proof}
Each factor lies in \([0,1]\), so the displayed values are nonnegative.
Summing over \(x_{n+1}\in\{0,1\}\) multiplies the value of
\(x_{1:n}\) by \(q+(1-q)=1\). Hence the laws are consistent, and since the
two words of length one have value \(1/2\) each, induction on \(n\) shows
that each sums to one. They define a finitely additive probability on the
finite-record algebra, so \cref{thm:p15-w02b-p15-022-finite-word-credences},
item~1, gives a unique law \(P_q\) with these finite laws; that item's proof
does not use the present theorem. Replacing every \(x_j\) by \(1-x_j\)
preserves each equality \(x_{j+1}=x_j\), so a word and its complement have
the same probability; hence \(P_q(X_j=1)=P_q(X_j=0)=1/2\).

(1) If \(q=1/2\), every word of length \(n\) has probability \(2^{-n}\),
which is rearrangement invariant. If \(q=1\), a word has probability
\(1/2\) when constant and zero otherwise; a rearrangement of a constant word
is constant and of a nonconstant word is nonconstant. Both laws are
exchangeable. If \(q\notin\{1/2,1\}\),
\[
 P_q(X_{1:3}=110)-P_q(X_{1:3}=101)
   =\tfrac12q(1-q)-\tfrac12(1-q)^2=\tfrac12(1-q)(2q-1)\neq0,
\]
while \(101\) is a rearrangement of \(110\). So \((X_1,X_2,X_3)\) is not
exchangeable.

(2) This contradicts item~1 of
\cref{thm:p15-w02b-p15-022-symmetry-consequences}.

(3) For \(0<q<1\),
\(P_q(0011)=\tfrac12q(1-q)q>0\) and \(P_q(00111)=\tfrac12q(1-q)q\cdot q\),
so the ratio is \(q\). Also \(P_q(1100)=\tfrac12q(1-q)q>0\) and
\(P_q(11001)=\tfrac12q(1-q)q(1-q)\), so the ratio is \(1-q\).
If \((X_1,\ldots,X_5)\) is exchangeable, so is \((X_1,\ldots,X_4)\), by
permutations fixing the fifth index. The words \(0011\) and \(1100\) are
rearrangements of each other, as are \(00111\) and \(11001\). The two
numerators and the two denominators are therefore equal, and so are the
predictive probabilities.
\end{proof}
\end{annexitem}

\begin{annexitem}{annex:p15-w02b-p15-022-finite-word-credences}{thm:p15-w02b-p15-022-finite-word-credences}{What coherent finite-record credences fix}
\begin{proof}
\(\mathcal A\) is an algebra: the complement of \(\{X_{1:n}\in B\}\) is
\(\{X_{1:n}\in B^c\}\), and for \(m\geq n\) the union with
\(\{X_{1:m}\in B'\}\) is
\(\{X_{1:m}\in(B\times\{0,1\}^{m-n})\cup B'\}\).

(1) For Borel \(D\subseteq\mathbb R^n\) put
\(\mu_n(D)=\sum_{b\in D\cap\{0,1\}^n}C(X_{1:n}=b)\). By finite additivity
of \(C\), \(\mu_n\) is a finite sum of point masses with total mass
\(C(\Omega)=1\), hence a Borel probability. Since
\(\{X_{1:n}=b\}\) is the disjoint union of \(\{X_{1:n+1}=b0\}\) and
\(\{X_{1:n+1}=b1\}\), \(\mu_{n+1}(D\times\mathbb R)=\mu_n(D)\); this is the
consistency of \cref{def:p15-consistent-marginals}.
\Cref{thm:p15-kolmogorov-countable} gives a probability \(\tilde P\) on
\(\mathbb R^{\mathbb N_+}\) with these finite laws. Let
\(h(t)=\mathbf1_{\{t=1\}}\) and \(H(\omega)=(h(\omega_j))_j\), which is
measurable into \(\Omega\) because each coordinate is. Put
\(P=\tilde P\circ H^{-1}\). For \(A=\{X_{1:n}\in B\}\),
\(H^{-1}(A)\) contains \(\pi_n^{-1}(B)\) and differs from it only inside
\(\pi_n^{-1}(\mathbb R^n\setminus\{0,1\}^n)\), which has
\(\tilde P\)-probability \(\mu_n(\mathbb R^n\setminus\{0,1\}^n)=0\). Hence
\(P(A)=\mu_n(B)=C(A)\). Uniqueness: \(\mathcal A\) is a pi-system
generating \(\mathcal F\), and two probabilities agreeing on it agree on
\(\mathcal F\) by the pi--lambda theorem.

The finite laws of \(P\) are the restrictions of \(C\) to words. The law of
\((X_1,\ldots,X_n)\) is invariant under permutations exactly when each word
has the value of each of its rearrangements, which is the stated condition.
The predictive formula concerns only the events \(\{X_{1:n}=b\}\) and
\(\{X_{1:n+1}=b1\}\) of \(\mathcal A\), on which \(C\) and \(P\) agree.

(2) Let \(A=\{X_{1:n}\in B\}\neq\varnothing\) and choose \(b\in B\).
The sequence \(b\) followed by ones has frequencies tending to one, so it lies
in \(A\cap L\). Construct \(w\) by following \(b\) with ones until the running
frequency exceeds \(2/3\), then zeros until it falls below \(1/3\), and so
on; each stage is finite because appending ones drives the frequency toward
one and appending zeros toward zero. The frequency of \(w\) has
\(\limsup\geq2/3\) and \(\liminf\leq1/3\), so \(w\in A\setminus L\).
Thus no nonempty member of \(\mathcal A\) lies in \(L\) or in \(L^c\). The
only member of \(\mathcal A\) inside \(L\) is \(\varnothing\), so the inner
bound is zero. If \(A\supseteq L\) with \(A\in\mathcal A\), then
\(A^c\in\mathcal A\) lies in \(L^c\), so \(A^c=\varnothing\) and the outer
bound is \(C(\Omega)=1\). For exchangeable \(P\),
\(P(L)=1\) by \cref{thm:p15-exchangeable-strong-laws}.

(3) A countably additive probability is finitely additive on
\(\mathcal A\). Non-exchangeability for \(q\notin\{1/2,1\}\) is
\cref{thm:p15-w02b-p15-022-markov-trend}, item~1.
\end{proof}
\end{annexitem}

\begin{annexitem}{annex:p15-w02b-p15-022-reflection}{thm:p15-w02b-p15-022-reflection}{Deference to a future conditioning self}
\begin{proof}
First, \(Q^N_m\) is a version of \(P(X_{N+1}=1\mid\mathcal F_m)\): it is
constant on each atom \(\{X_{1:m}=b\}\), so \(\mathcal F_m\)-measurable,
and each \(D\in\mathcal F_m\) is a union of atoms, whence
\[
 P(\{X_{N+1}=1\}\cap D)
   =\sum_{b}Q^N_m(b)\,P(X_{1:m}=b)
   =E(Q^N_m;D),
\]
where the sum runs over the positive-probability atoms \(\{X_{1:m}=b\}\)
contained in \(D\), \(Q^N_m(b)\) is the ratio on that atom, and null atoms
contribute zero to both sides.

(1) By \cref{thm:p15-w02b-p15-022-predictive}, item~1,
\(Q^N_m=E(\Theta\mid\mathcal F_m)\) almost surely. For \(n\leq m\),
\(\mathcal F_n\subseteq\mathcal F_m\), and the tower property gives
\(E(Q^N_m\mid\mathcal F_n)=E(\Theta\mid\mathcal F_n)=Q^N_n\). Also
\(Q^N_m\) and \(Q^m_m\) are both versions of \(E(\Theta\mid\mathcal F_m)\)
constant on atoms; they agree on every positive-probability atom and are
both zero on the null atoms.

(2) The event \(D=\{X_{1:n}=c\}\cap\{Q^N_m=x\}\) lies in \(\mathcal F_m\)
because \(n\leq m\). By the version property,
\(P(\{X_{N+1}=1\}\cap D)=E(Q^N_m;D)=xP(D)\). Divide by \(P(D)>0\).

(3) By \cref{thm:p15-w02b-p15-022-mixing-law}, item~2, \(P'\) gives
\(P'(X_2=1\mid X_1=1)=2/3\) and \(P'(X_2=1\mid X_1=0)=1/3\), so
\(\{R_1=2/3\}=\{X_1=1\}\). Under the independent fair law,
\(P(X_2=1\mid X_1=1)=1/2\).
\end{proof}
\end{annexitem}

\begin{annexitem}{annex:p15-w02b-p15-022-reverse-strong-law-boundary}{thm:p15-w02b-p15-022-reverse-strong-law-boundary}{Where the reverse strong law stops}
\begin{proof}
(1) For \(V\geq0\), \(V\leq\sum_{k\geq0}\mathbf1_{\{V>k\}}\): if
\(m\leq V<m+1\) for an integer \(m\geq0\), the sum counts at least
\(k=0,\ldots,m-1\), and also \(k=m\) when \(V>m\). Taking expectations by
monotone convergence and using identical laws,
\[
 \infty=E|X_1|\leq1+\sum_{n\geq1}P(|X_1|>n)=1+\sum_{n\geq1}P(|X_n|>n).
\]
The events \(\{|X_n|>n\}\) are mutually independent, so
\cref{thm:p15-borel-cantelli-second} gives \(P(|X_n|>n\text{ i.o.})=1\).
Let \(G\) be the measurable event on which \(S_n/n\) converges to a real
number \(c\). On \(G\), for \(n\geq2\),
\[
 \frac{X_n}{n}=\frac{S_n}{n}-\frac{n-1}{n}\cdot\frac{S_{n-1}}{n-1}
   \longrightarrow c-c=0,
\]
which is incompatible with \(|X_n|/n>1\) infinitely often. Hence
\(G\) is contained in a null set.

(2) Let \(\pi\) permute \(1,\ldots,n\). Then \(S_n\circ T_\pi=S_n\) and
\(X_j\circ T_\pi=X_j\) for \(j>n\), so the generating events
\(\{S_n\in B\}\) and \(\{X_j\in B\}\) of \(\mathcal H_n\) are invariant
and lie in the sigma-field \(\mathcal E_n\)
(\cref{thm:p15-exchangeable-canonical}); hence
\(\mathcal H_n\subseteq\mathcal E_n\). By
\cref{thm:p15-exchangeable-averages} with \(k=1\) and \(\varphi(x)=x\),
\(S_n/n=E(X_1\mid\mathcal E_n)\). Since \(S_n/n\) is
\(\mathcal H_n\)-measurable, the tower property gives
\(E(X_1\mid\mathcal H_n)=E(E(X_1\mid\mathcal E_n)\mid\mathcal H_n)=S_n/n\).
The fields decrease: \(S_{n+1}=S_n+X_{n+1}\) and \(X_j\), \(j>n+1\), are
\(\mathcal H_n\)-measurable. \Cref{thm:p15-levy-backward} gives the
almost-sure and \(L^1\) limit \(E(X_1\mid\mathcal H_\infty)\).

(3) The law of \((X_1,\ldots,X_n)\) is that of \((\xi,\ldots,\xi)\), which is
unchanged by any permutation of identical coordinates; so the law is
exchangeable, and \(P(X_j\neq X_1)=0\), whence \(S_n/n=X_1\) almost surely.
Let \(F\) be the distribution function of \(\xi\). If \(F\) took only the
values \(0\) and \(1\), then with \(c^*=\inf\{c:F(c)=1\}\), finite because
\(F\) tends to \(0\) and \(1\) at \(\mp\infty\), right continuity would give
\(F(c^*)=1\) and \(F(c)=0\) for \(c<c^*\), so \(\xi=c^*\) almost surely.
Hence some \(c\) has \(0<F(c)<1\). Put \(\bar L=\limsup_nS_n/n\). For each
\(m\), \(\bar L\) is a function of \(S_n\), \(n\geq m\), all
\(\mathcal H_m\)-measurable; so \(\{\bar L\leq c\}\in\mathcal H_\infty\),
and \(P(\bar L\leq c)=P(X_1\leq c)=F(c)\in(0,1)\). The limit \(X_1\)
is not almost surely equal to \(E\xi\), since \(\xi\) is not almost surely
constant.
\end{proof}
\end{annexitem}

\begin{annexitem}{annex:p15-w02b-p15-022-backward-changing-boundary}{thm:p15-w02b-p15-022-backward-changing-boundary}{The envelope in the change-of-integrand theorem}
\begin{proof}
(1) The sequence \(\mathcal G_n=\mathcal G=\sigma(B_1,B_2,\ldots)\)
satisfies \(\mathcal G_{n+1}\subseteq\mathcal G_n\) and
\(\bigcap_n\mathcal G_n=\mathcal G\). \Cref{thm:p15-ui-counterexamples},
item~3, proves that \(Z_n=nB_nC_n\to0\) almost surely, that \((Z_n)\) is
uniformly integrable, and that \(E(Z_n\mid\mathcal G)=B_n\) fails to
converge almost surely. If an integrable \(W\geq|Z_n|\) existed for all
\(n\), \cref{thm:p15-backward-changing-variables} would give almost-sure
convergence of \(E(Z_n\mid\mathcal G_n)\); so no such \(W\) exists.

(2) For \(\omega\in(0,1)\) and \(n>1/\omega\), \(Z_n(\omega)=0\). Each
\(EZ_n=n\cdot(1/n)=1\), and conditional expectation given the trivial field
is the mean. The limit \(Z=0\) has \(E(Z\mid\mathcal G_\infty)=0\).

(3) For integrable \(V\), \(-|V|\leq V\leq|V|\) and monotonicity give
\(|E(V\mid\mathcal G)|\leq E(|V|\mid\mathcal G)\), so
\(E|E(V\mid\mathcal G)|\leq E|V|\). Hence
\[
 E\bigl|E(Z_n\mid\mathcal G_n)-E(Z\mid\mathcal G_\infty)\bigr|
   \leq E|Z_n-Z|+E\bigl|E(Z\mid\mathcal G_n)-E(Z\mid\mathcal G_\infty)\bigr|.
\]
The first term tends to zero by hypothesis, and the second by
\cref{thm:p15-levy-backward}. In item~1, \(E|B_n|=1/n\to0\).
\end{proof}
\end{annexitem}

\begin{annexitem}{annex:p15-w02b-p15-022-projection-boundary}{thm:p15-w02b-p15-022-projection-boundary}{The independence hypothesis in the conditional-projection lemma}
\begin{proof}
(1) \(V\) is \(\sigma(V)\)-measurable and integrable, and it has the same
integral as itself over every \(A\in\sigma(V)\); so it satisfies the defining
property of \(E(V\mid\sigma(V))\) and is a version of \(Y\). It is
\(\mathcal H\)-measurable because \(\mathcal H=\sigma(V)\). Let \(F\) be the
distribution function of \(V\). If \(F\) took only the values \(0\) and
\(1\), then with \(c^*=\inf\{c:F(c)=1\}\), finite because \(F\) tends to
\(0\) and \(1\) at \(\mp\infty\), right continuity would give
\(F(c^*)=1\) and \(F(c)=0\) for \(c<c^*\), so \(V=c^*\) almost surely,
contrary to hypothesis. Hence some \(c\) has \(0<F(c)<1\). The event
\(A=\{V\leq c\}\) lies in \(\mathcal H\) and in \(\sigma(V)\), and
\[
 P(A\cap\{V\leq c\})=P(A)=F(c)\neq F(c)^2=P(A)\,P(V\leq c),
\]
so \(\mathcal H\) is not independent of \(V\). If \(Y=EV\) almost surely,
then \(V=Y=EV\) almost surely, again contrary to hypothesis; so
\(P(Y\neq EV)>0\). For the fair binary \(V\), \(\{V=1\}=\{V>1/2\}\), and
the displayed numbers follow from \(P(V=1)=1/2\).

(2) The law of \(X\) is \(\tfrac12\delta_{(0,0,\ldots)}+\tfrac12\delta_{(1,1,\ldots)}\).
The vector \((X_1,\ldots,X_n)\) equals \((\xi,\ldots,\xi)\), whose law is
unchanged by any permutation of its identical coordinates; so the law is
exchangeable. Put \(\bar Y=\limsup_nS_n/n\), with values in \([0,1]\).

\emph{\(\mathcal E\)-measurability.} Fix \(n\) and a permutation \(\pi\) of
\(1,\ldots,n\). For \(m\geq n\), \(S_m\circ T_\pi=S_m\), because \(T_\pi\)
only rearranges the first \(n\) coordinates, all of which enter \(S_m\).
Since \(\bar Y=\limsup_{m\geq n}S_m/m\), also \(\bar Y\circ T_\pi=\bar Y\),
so \(T_\pi^{-1}\{\bar Y\leq c\}=\{\bar Y\leq c\}\) for every real \(c\). Thus
\(\bar Y\) is \(\mathcal E_n\)-measurable for every \(n\), hence
\(\mathcal E\)-measurable.

\emph{\(\mathcal H\)-measurability.} \(S_m/m-(S_m-X_1)/m=X_1/m\to0\)
everywhere, so \(\bar Y=\limsup_m(S_m-X_1)/m\), a function of
\(X_2,X_3,\ldots\).

\emph{Version.} Almost surely the outcome is a constant sequence, on which
\(S_m/m=X_1\) for every \(m\); so \(\bar Y=X_1\) almost surely. For
\(A\in\mathcal E\), therefore, \(E(X_1;A)=E(\bar Y;A)\), and \(\bar Y\) is
\(\mathcal E\)-measurable and integrable; it is a version of
\(E(X_1\mid\mathcal E)\).

\emph{Failure.} \(P(X_1=1,X_2=1)=P(\xi=1)=1/2\), whereas
\(P(X_1=1)P(X_2=1)=1/4\), so \(\mathcal H\ni\{X_2=1\}\) is not independent
of \(V=X_1\). Finally \(EV=1/2\), while \(\bar Y=X_1\in\{0,1\}\) almost
surely.

(3) Let \(A=\{Y'>EV\}\in\mathcal H\) and \(A_0=\{Y>EV\}\in\mathcal G\),
where \(Y\) is the given \(\mathcal G\)-measurable version. Since
\(Y=Y'\) almost surely, \(A\) and \(A_0\) differ by a null set, and
\[
 \begin{aligned}
 E(Y'-EV;A)&=E(Y;A_0)-EV\cdot P(A_0)\\
   &=E(V;A_0)-EV\cdot P(A)\\
   &=E(V;A)-EV\cdot P(A)=0,
 \end{aligned}
\]
using the defining property of \(Y\) on \(A_0\in\mathcal G\) and then the
hypothesis on \(A\in\mathcal H\). The integrand \(Y'-EV\) is strictly
positive on \(A\), so \(P(A)=0\). The same argument with
\(\{Y'<EV\}\) gives \(P(Y'<EV)=0\). Hence \(Y'=EV\), and so \(Y=EV\),
almost surely. Only integrability of \(V\) has been used.

If \(\mathcal H\) is independent of \(V\) and \(A\in\mathcal H\), then
\(V\) and \(\mathbf1_A\) are independent and integrable, so
\(E(V\mathbf1_A)=EV\cdot P(A)\); the hypothesis of item~3 holds.

For the example, \(EV=E\varepsilon\,EW=0\) by independence. Every
\(A\in\sigma(W)\) has the form \(\{W\in B\}\), and \(\varepsilon\) is
independent of \(W\mathbf1_{\{W\in B\}}\), so
\(E(V;A)=E\varepsilon\cdot E(W\mathbf1_{\{W\in B\}})=0=EV\cdot P(A)\). This
is the hypothesis of item~3 for \(\mathcal H=\sigma(W)\), and it also shows
that the constant \(0\) is a version of \(E(V\mid\sigma(W))\). Finally,
\(|V|=W\), so \(\{|V|=1\}=\{W=1\}\) and
\(P(W=1,|V|=1)=1/2\), while \(P(W=1)P(|V|=1)=1/4\). Thus \(\sigma(W)\),
which is both \(\mathcal G\) and \(\mathcal H\), is not independent of
\(V\).
\end{proof}
\end{annexitem}

\begin{annexitem}{annex:p15-w02b-p15-022-presentable}{thm:p15-w02b-p15-022-presentable}{Where the regular conditional law enters}
\begin{proof}
Two measurability facts are used throughout. For fixed
\(B_1,\ldots,B_n\in\mathcal S\), the map
\(\mu\mapsto\prod_{j=1}^n\mu(B_j)=\prod_{j=1}^ne_{B_j}(\mu)\) is a finite
product of generating evaluations, hence measurable and bounded by \(1\)
on \(\mathscr P(S)\). And two probability measures on
\(\mathcal S^{\otimes n}\) that agree on the measurable rectangles
\(B_1\times\cdots\times B_n\) agree on \(\mathcal S^{\otimes n}\), because
the rectangles form a pi-system generating \(\mathcal S^{\otimes n}\) and
the class of sets on which two probability measures agree is a
lambda-system.

Item~1. Let \(P\) be presentable with mixing law \(\Lambda\), let
\(n\geq1\), and let \(\pi\) be a permutation of \(\{1,\ldots,n\}\). The law
of \((X_{\pi(1)},\ldots,X_{\pi(n)})\) gives the rectangle
\(B_1\times\cdots\times B_n\) the mass
\[
 P(X_{\pi(1)}\in B_1,\ldots,X_{\pi(n)}\in B_n)
   =P\bigl(X_1\in B_{\pi^{-1}(1)},\ldots,X_n\in B_{\pi^{-1}(n)}\bigr)
   =\int\prod_{i=1}^n\mu\bigl(B_{\pi^{-1}(i)}\bigr)\,\Lambda(d\mu).
\]
The first equality only reorders the conditions, since
\(X_{\pi(j)}\in B_j\) for all \(j\) says that \(X_i\in B_{\pi^{-1}(i)}\)
for all \(i\). The second is presentability. The integrand is a product
of the same \(n\) real factors \(\mu(B_1),\ldots,\mu(B_n)\) in another
order, so it equals \(\prod_{j=1}^n\mu(B_j)\), and the integral equals
\(P(X_1\in B_1,\ldots,X_n\in B_n)\). The laws of
\((X_{\pi(1)},\ldots,X_{\pi(n)})\) and \((X_1,\ldots,X_n)\) therefore agree
on rectangles and hence on \(\mathcal S^{\otimes n}\). As \(n\) and
\(\pi\) were arbitrary, \(P\) is exchangeable. No property of \(S\) beyond
being a measurable space was used.

Item~2. The sigma-field of \(\mathscr P(S)\) is generated by the maps
\(e_B\), and \(e_B\circ K=K(B)\) is \(\mathcal E\)-measurable for every
\(B\in\mathcal S\). Hence \(K\) is measurable from
\((\Omega,\mathcal E)\) into \(\mathscr P(S)\), and so also from
\((\Omega,\mathcal F)\); its image law \(\Lambda=P\circ K^{-1}\) is a
probability on \(\mathscr P(S)\). Apply
\cref{thm:p15-definetti-factorization}, which holds on an arbitrary
measurable coordinate space, with the bounded measurable functions
\(f_j=\mathbf1_{B_j}\):
\[
 P(X_1\in B_1,\ldots,X_n\in B_n\mid\mathcal E)
   =E\Bigl(\prod_{j=1}^n\mathbf1_{B_j}(X_j)\Bigm|\mathcal E\Bigr)
   =\prod_{j=1}^nE\bigl(\mathbf1_{B_j}(X_1)\bigm|\mathcal E\bigr)
   =\prod_{j=1}^nK(B_j)
\]
almost surely. The last equality holds because \(K(B_j)\) is a version of
\(P(X_1\in B_j\mid\mathcal E)\) and a finite product of almost surely
equal factors is almost surely equal. Taking expectations, and using
\(\prod_jK(B_j)=\bigl(\prod_je_{B_j}\bigr)\circ K\) with the change of
variables for the image law,
\[
 P(X_1\in B_1,\ldots,X_n\in B_n)
   =E\prod_{j=1}^nK(B_j)
   =\int_{\mathscr P(S)}\prod_{j=1}^n\mu(B_j)\,\Lambda(d\mu).
\]
This is presentability with the stated \(\Lambda\).

Item~3. For standard Borel \(S\), a regular conditional law of \(X_1\)
given \(\mathcal E\) exists by the construction of
\cref{sec:p15-regular-conditional-laws}, which is the choice made in
\cref{thm:p15-definetti-kernel}; that theorem's evaluation sigma-field is
the one used here. Item~2 then gives presentability. The cylinder formula
of \cref{thm:p15-definetti-kernel} is the same conclusion.

Item~4. If an exchangeable law admitted a regular conditional law of
\(X_1\) given \(\mathcal E\), item~2 would make it presentable.
\end{proof}
\end{annexitem}

\begin{annexitem}{annex:p15-w02b-p15-022-alternatives-contrast}{thm:p15-w02b-p15-022-alternatives-contrast}{One state, one piece of evidence, five updates}
\begin{proof}
Write \(\pi_j=\pi(w_j)\), \(\pi_4=0\), and \(F_i=\{w_1,\ldots,w_i\}\), so
that \(m(F_i)=\pi_i-\pi_{i+1}\): \(1-\tfrac12=\tfrac12\),
\(\tfrac12-\tfrac14=\tfrac14\), \(\tfrac14-0=\tfrac14\). The masses are
nonnegative and sum to one.

(1) The focal sets containing \(w_j\) are \(F_j,\ldots,F_3\), so
\(\operatorname{Pl}_m(\{w_j\})=\sum_{i\geq j}(\pi_i-\pi_{i+1})=\pi_j\).
For \(A=W\), \(\operatorname{Bel}_m(W)=1\) and
\(N_\kappa(W)=1-2^{-\kappa(\varnothing)}=1\). Otherwise let \(j\) be the
least index with \(w_j\notin A\). Then \(F_i\subseteq A\) exactly when
\(i<j\), and
\[
 \operatorname{Bel}_m(A)=\sum_{i<j}(\pi_i-\pi_{i+1})=\pi_1-\pi_j=1-\pi_j .
\]
Since \(t\mapsto2^{-t}\) is decreasing,
\(2^{-\kappa(A^c)}=2^{-\min_{w\notin A}\kappa(w)}=\max_{w\notin A}\pi(w)\),
and because \(\pi_1>\pi_2>\pi_3\) this maximum is \(\pi_j\). Hence
\(N_\kappa(A)=1-\pi_j=\operatorname{Bel}_m(A)\).
By the pignistic formula of \cref{thm:p15-r47-pignistic},
\[
 p(w_1)=\tfrac12+\tfrac18+\tfrac1{12}=\tfrac{17}{24},\qquad
 p(w_2)=\tfrac18+\tfrac1{12}=\tfrac5{24},\qquad
 p(w_3)=\tfrac1{12}.
\]
At threshold zero, \(A\) is believed when \(\kappa(\neg A)>0\), that is when
no world of rank zero lies outside \(A\), that is when \(w_1\in A\).
Likewise \(N_\kappa(A)>0\) exactly when \(\kappa(A^c)>0\), and \(w_1\) is the
unique world of minimal rank.

(2a) \(p(E)=\tfrac5{24}+\tfrac2{24}=\tfrac7{24}>0\), so
\(p(w_2\mid E)=\tfrac57\), \(p(w_3\mid E)=\tfrac27\) and
\(p(w_1\mid E)=0\).

(2b) \(\kappa(E)=\min\{1,2\}=1<\infty\), so by
\cref{def:p15-rat53-conditionals}
\(\kappa(\{w\}\mid E)=\kappa(\{w\}\cap E)-1\): this is
\(\kappa(\varnothing)-1=\infty\) for \(w_1\), \(0\) for \(w_2\) and \(1\) for
\(w_3\). As in item~1, \(A\) is believed at threshold zero after the update
exactly when the unique rank-zero world \(w_2\) lies in \(A\).

(2c) The \(E\)-worlds are \(w_2,w_3\), and \(w_2\prec w_3\).

(2d) \(\Pi(E)=\max\{\tfrac12,\tfrac14\}=\tfrac12\), so
\(\pi_E(w_2)=1\) and \(\pi_E(w_3)=\tfrac12\); and
\(2^{-0}=1\), \(2^{-1}=\tfrac12\), \(2^{-\infty}=0\).

(2e) Combine \(m\) with the assignment of mass one to \(E\). The products
fall on \(\{w_1\}\cap E=\varnothing\) with weight \(\tfrac12\), on
\(F_2\cap E=\{w_2\}\) with weight \(\tfrac14\), and on \(W\cap E=E\) with
weight \(\tfrac14\). So \(K=\tfrac14+\tfrac14=\tfrac12>0\), and
\(m_E(\{w_2\})=m_E(E)=\tfrac{1/4}{1/2}=\tfrac12\). Both focal sets contain
\(w_2\); only \(E\) contains \(w_3\); neither contains \(w_1\). So the
singleton plausibilities are \(0,1,\tfrac12\), which is \(\pi_E\); and
\(\operatorname{Bel}_{m_E}(\{w_2\})=m_E(\{w_2\})=\tfrac12\). The pignistic
probability of \(m_E\) is \(\tfrac12+\tfrac14=\tfrac34\) at \(w_2\) and
\(\tfrac14\) at \(w_3\); and \(\tfrac34\neq\tfrac57\).

(2f) Obtain \(q^a\) by giving the mass of \(\{w_1\}\) and of \(F_2\) to
\(w_1\) and the mass of \(W\) to \(w_3\); obtain \(q^b\) the same way but
giving the mass of \(W\) to \(w_2\). In each case every focal set sends its
mass to one of its own elements, so for every \(A\) each focal
\(F\subseteq A\) contributes its whole mass to \(A\), and
\(q(A)\geq\sum_{F\subseteq A}m(F)=\operatorname{Bel}_m(A)\). Both give
\(E\) the mass \(\tfrac14\) of \(W\); hence \(q^a(\{w_2\}\mid E)=0\) and
\(q^b(\{w_2\}\mid E)=\tfrac{1/4}{1/4}=1\). For \(0\leq t\leq1\), the
mixture \(q_t=(1-t)q^a+tq^b\) again dominates \(\operatorname{Bel}_m\),
has \(q_t(E)=\tfrac14\), and has \(q_t(\{w_2\}\mid E)=t\). Every
conditional probability lies in \([0,1]\), so the set of values
\(q(\{w_2\}\mid E)\), over dominating \(q\) with \(q(E)>0\), is exactly
\([0,1]\). By~(e),
\(\operatorname{Bel}_{m_E}(\{w_2\})=\tfrac12\) and
\(\operatorname{Pl}_{m_E}(\{w_2\})=1\).

(3) After~(a), \(p(E^c\mid E)=0\), so the ratio conditional given \(E^c\)
has a zero denominator, and \(\{E,E^c\}\) has a cell of zero old weight,
which \cref{def:p15-rat41-jeffrey} excludes. After~(b),
\(\kappa(\neg E\mid E)=\infty\), whereas \cref{def:p15-rat53-soft-update}
requires both cells of the input to have finite rank. For the firmness
updates, \(\kappa(E)=1\) and \(\kappa(\neg E)=0\) are finite; with
\(a=0\), \(b=1\) the world-rank formula of \cref{thm:p15-rat53-update}
gives \(w_2\mapsto1-1+0=0\),
\(w_3\mapsto2-1+0=1\) and \(w_1\mapsto0-0+1=1\). In the new ranking
\(\{w_1\}\) has rank \(1\) and its complement rank \(0\); accepting
\(\{w_1\}\) with \(a=0\), \(b=1\) gives \(w_1\mapsto1-1+0=0\),
\(w_2\mapsto0-0+1=1\) and \(w_3\mapsto1-0+1=2\). Minimal-world selection
from the one-element set \(\{w_1\}\) returns it in every order.

(4) Selling a unit bet on \(A\) at price \(x\) yields \(x-\mathbf1_A\).
Here \(\operatorname{Bel}_m(\{w_1\})=\tfrac12\) and
\(\operatorname{Bel}_m(E)=0\) by item~1. Selling both bets yields
\(\tfrac12-\mathbf1_{\{w_1\}}-\mathbf1_E=\tfrac12-1=-\tfrac12\) in every
world, since \(\{w_1\}\) and \(E\) partition \(W\).

(5) By \cref{def:p15-r47-ranking},
\(N_\kappa(\varnothing)=1-2^{-\kappa(W)}=0\) and
\(N_\kappa(\{w_2\})=1-2^{-\kappa(\{w_1,w_3\})}=1-2^{0}=0\), while
\(N_\kappa(W)=1-2^{-\infty}=1\) and
\(N_\kappa(\{w_1,w_3\})=1-2^{-\kappa(\{w_2\})}=\tfrac12\). A function
\(n\) with \(N_\kappa(A^c)=n(N_\kappa(A))\) would need \(n(0)=1\) and
\(n(0)=\tfrac12\). By item~1 the same values hold for
\(\operatorname{Bel}_m\). For a probability \(q\), finite additivity gives
\(q(A)+q(A^c)=q(W)=1\).
\end{proof}
\end{annexitem}
